\iffull
Four of this paper's statistics could not fail. A recovery curve in which the
estimator cancels still draws a smooth line; a grid that cannot resolve a
crossing still returns a stable number; a confound diagnostic whose tolerance
sits above its own attainable bound still reports clean; and a premise check
reading a control detected in 99\% of cells reports comparable arms while those
same controls shift by up to $0.71$ on log$_2$ expression. Each looked like a
measurement, and in every case the fix was to ask what the statistic was a
function of rather than to look harder at its output. \textbf{We wrote the
fourth ourselves, to this rule, while writing this paper}, and it took a
bootstrap interval to notice --- which is the strongest claim we can make that
the failure is structural rather than careless.
\fi

Asking what a statistic is a function of has a cheap answer for recovery curves
specifically, and the answer generalises past our instance: compare against the
realised draw rather than the requested parameter, and assert the difference is
not identically zero.
It costs one line, it runs on the sweep you already have, and in a patient
simulator it separates two estimators whose curves agree to machine precision.
We would rather it had been in our design document than in our findings.

\iffull\else
A synthetic benchmark (Appendix~\ref{app:benchtable}) locates what that rule was
buying: three ungated methods invent an effect in $200/200$ replicates where the
estimand does not exist, but a width gate suppresses that as completely as our
count gate while costing no detections. The finiteness check, not the threshold,
does the work.
\fi
